\documentclass[11pt]{article}

\usepackage[margin=1in]{geometry}
\usepackage{booktabs}
\usepackage{amsmath}
\usepackage{hyperref}

\title{Dual-Domain MMD Adaptation: Experiment Summary}
\author{}
\date{\today}

\begin{document}
\maketitle

\section{What was built}

A YOLOv10n detector with an added Maximum Mean Discrepancy (MMD) alignment term between two
domains --- \textbf{real} (VisDrone aerial imagery) and \textbf{syn} (Syndrone synthetic aerial
imagery), single class (vehicle). The pipeline follows a two-stage procedure:

\begin{enumerate}
    \item \textbf{Pretraining} (standard, stock ultralytics, no MMD): the model is trained
    normally on one domain, called the \emph{source} domain.
    \item \textbf{Adaptation}: that checkpoint is then further fine-tuned on the \emph{other}
    domain, called the \emph{target} domain, using the target domain's own detection loss plus an
    MMD term computed between the two domains' backbone features (hooked at layer 10, RBF
    kernel, momentum-smoothed bandwidth, constant weight $\lambda_{DA}=0.8$). The source domain's
    forward pass is gradient-detached (\texttt{detach\_source\_features=True}), so it only ever
    supplies a fixed reference feature distribution --- it never receives a training signal
    itself during adaptation; only the target domain's features are pulled toward it.
\end{enumerate}

A deliberate size asymmetry was built into the setup: the \emph{source} domain always uses a
small, stride-3 subsampled dataset ($\sim\!1/3$ of full size --- real: 2157 train images, syn:
1509), while the \emph{target} domain always uses the \emph{full} dataset (real: 6470, syn: 4525)
--- simulating ``pretrained on limited data, fine-tune using a larger dataset from a different
domain.''

\section{What was tested, and in which direction}

\textbf{Baseline (no adaptation)}: each domain pretrained standalone on its own small
subsampled data, then cross-evaluated on both domains' (small) validation sets --- giving
in-domain and zero-shot cross-domain reference numbers.

\textbf{Direction 1 --- \texttt{real\_source\_syn\_target}}: pretrained on \textbf{real} (small),
adapted with \textbf{syn} as the target (full dataset). Tests whether real$\to$synthetic
adaptation helps synthetic-domain performance.

\textbf{Direction 2 --- \texttt{syn\_source\_real\_target}}: pretrained on \textbf{syn} (small),
adapted with \textbf{real} as the target (full dataset). Tests whether synthetic$\to$real
adaptation helps real-domain performance --- closer to the project's original motivating premise
(cheap synthetic data helping a scarce-label real domain).

\section{Results}

\begin{table}[h!]
    \centering
    \begin{tabular}{lrr}
        \toprule
         & mAP50 & mAP50-95 \\
        \midrule
        baseline syn$\to$syn (in-domain) & 0.758 & 0.508 \\
        baseline real$\to$real (in-domain) & 0.896 & 0.641 \\
        baseline syn$\to$real (cross, no adaptation) & 0.194 & 0.116 \\
        baseline real$\to$syn (cross, no adaptation) & 0.592 & 0.359 \\
        \midrule
        \textbf{Dir. 1} --- adapted, eval on \textbf{syn (target)} & \textbf{0.824} (+0.067) & \textbf{0.600} (+0.091) \\
        Dir. 1 --- adapted, eval on real (source) & 0.489 ($-$0.407) & 0.340 ($-$0.301) \\
        \midrule
        \textbf{Dir. 2} --- adapted, eval on \textbf{real (target)} & \textbf{0.867} ($-$0.029) & \textbf{0.599} ($-$0.042) \\
        Dir. 2 --- adapted, eval on syn (source) & 0.568 ($-$0.190) & 0.337 ($-$0.171) \\
        \bottomrule
    \end{tabular}
    \caption{Baseline vs.\ adapted mAP, both directions. Deltas are relative to the matching
    in-domain baseline.}
\end{table}

\textbf{Direction 1 (target = syn) genuinely improved over its baseline} on both metrics.
\textbf{Direction 2 (target = real) did not} --- a small regression at both IoU thresholds. In
both directions, the \textbf{source domain's own performance dropped substantially} ---
the source never receives a detection-loss gradient for the full 100-epoch adaptation run, so the
shared backbone drifts away from what worked for it, a forgetting effect independent of whether
MMD is ``working.''

\section{Feature-space (PCA) analysis}

A fixed-basis PCA (fit once before adaptation on 500 GAP-pooled layer-10 features per domain,
frozen thereafter, re-projected every epoch) was tracked for Direction 1. It showed the target
(syn) domain's features collapsing into a tight cluster near one edge of the source (real)
domain's spread, rather than spreading out to match the source distribution's shape ---
consistent with a known characteristic of naive two-sample MMD alignment (collapsing toward a
nearby point reduces the MMD statistic more cheaply than reshaping variance to genuinely match
the source distribution). The source domain's own spread also contracted somewhat over training,
despite never receiving gradient directly --- a side effect of the shared backbone being updated
by the target domain's loss.

\section{Not yet tested (planned follow-ups)}

\begin{itemize}
    \item \textbf{\texttt{detach\_source\_features=False}} (mutual pull): letting both domains'
    features move toward each other, to see whether it avoids the source-domain regression or
    produces closer convergence.
    \item \textbf{Detection loss on both domains simultaneously} (plus MMD): a fuller
    joint-objective variant, not yet implemented in the model.
    \item \textbf{Reversed size pairing}: pretrain on the \emph{full} dataset (source), adapt
    into the \emph{small} subsampled dataset (target) --- the mirror image of the current setup,
    closer to ``lots of cheap data helps a domain with scarce data.''
\end{itemize}

\end{document}
